% ============================================================================
% Fichier  : partie5/P05_C28_auditer_ia.tex
% Titre    : Auditer l'Intelligence Artificielle, Auditer avec Elle
% Auteur   : MINKA MI NGUIDJOÏ Thierry Emmanuel
% Version  : 1.0 -- novembre 2026
% Base   : ACSI C20, repris et développé pour INF3226
% ============================================================================
% Sources : cours GCI (MSN-2S3-MSI03), chapitre 19 (gouvernance de l'IA,
% six risques, ratio d'impact et seuil de 0,8 présenté comme signal sans
% valeur juridique au Cameroun) ; ISO/IEC 42001:2023 ; NIST AI RMF 1.0
% (quatre fonctions) ; Recommandation de l'UNESCO (2021) ; article A0
% (hal-05701007), section « Le défi de l'IA générative » : outil transversal
% d'opérationnalisation des piliers, non sixième pilier, position
% révisable. La section « preuve judiciaire » d'ACSI C20 est développée au
% Chapitre 29. Modèle de notation des microcrédits de la
% BANQUE SANAGA : construit pour ce chapitre.
% ============================================================================

\chapter{Auditer l'Intelligence Artificielle, Auditer avec Elle}
\label{chap:audit-ia}
\chaptermark{Auditer l'IA}

\epigraph{%
    « Un modèle qui décide doit pouvoir être audité.\\
    Un modèle qui aide l'auditeur ne décide de rien. »%
}{Règle du cours}

\niveauChapitre{Expert}{%
    Maîtrise du Chapitre~\ref{chap:audit-continu} (analyse de données) et
    du Chapitre~\ref{chap:audit-projets} (changements, mise en production)%
}

\begin{boxobjectifs}
\begin{itemize}
    \item Identifier, du point de vue de l'auditeur, les risques d'un
          système d'IA~: biais, explicabilité, dérive.
    \item Construire le programme d'audit d'un modèle, et mesurer un biais.
    \item Évaluer l'apport et les limites de l'IA générative dans la
          mission, et fixer ses règles d'usage.
    \item Situer les cadres de l'IA~: ISO/IEC~42001, NIST AI RMF, et le
          règlement européen sur l'IA comme exemple de régulation.
\end{itemize}
\end{boxobjectifs}

\bigskip

% ============================================================================
\section*{Question directrice : qui audite la machine, et avec quoi~?}
% ============================================================================

Depuis 2025, les agences de la BANQUE SANAGA accordent les microcrédits à
l'aide d'un modèle de notation fourni par un prestataire. Le comité
d'audit s'inquiète~: les refus semblent plus fréquents dans certaines
régions, et le taux de défaut augmente. Il demande à Mme~EKAMBI d'auditer
le modèle.

Chez LiZbiZ, un auditeur junior d'AUDIX CONSEIL a résumé les quarante
procédures de l'entreprise avec un assistant d'IA générative.

\medskip
\textbf{Plan du chapitre.}
\begin{enumerate}
    \item Ce que l'IA change pour l'auditeur (section~\ref{sec:c28-risques}).
    \item Le programme d'audit d'un modèle (section~\ref{sec:c28-programme}).
    \item Auditer avec l'IA générative (section~\ref{sec:c28-avec}).
\end{enumerate}

\separateur

% ============================================================================
\section{Ce que l'IA Change pour l'Auditeur}
\label{sec:c28-risques}
% ============================================================================

Le cours GCI (MSN-2S3-MSI03, chapitre~19) présente la gouvernance de
l'IA~: ses trois traits (un système qui apprend des données, s'explique
mal et évolue), ses six risques et ses cadres de référence. L'auditeur en
retient trois conséquences.

\begin{itemize}
    \item \textbf{Les règles ne sont pas écrites}~: dans une application
          classique, le critère est dans le code (un seuil, un calcul). Dans
          un modèle, il est appris. L'auditeur ne lit pas la règle~; il
          observe les décisions.
    \item \textbf{Le contrôle se déplace vers les données et la
          surveillance}~: provenance et représentativité des données
          d'entraînement, mesure des écarts entre groupes, suivi de la
          performance en production.
    \item \textbf{Chaque réentraînement est un changement}~: un modèle
          réentraîné est un nouveau modèle, qui relève de la gestion des
          changements et d'une nouvelle validation
          (Chapitre~\ref{chap:audit-projets}).
\end{itemize}

\begin{tableaucours}{L{3.6cm}L{2.0cm}L{6.6cm}}{%
    \tableauHeader{Référence} &
    \tableauHeader{Nature} &
    \tableauHeader{Usage pour l'auditeur}
}
ISO/IEC 42001 (2023) & Norme certifiable & Exigences d'un système de
management de l'IA~: critères pour un audit de type SMSI
(Chapitre~\ref{chap:audit-technique})~\cite{iso42001_2023} \\
NIST AI RMF 1.0 (2023) & Cadre & Quatre fonctions~: gouverner,
cartographier, mesurer, gérer~; trame de décomposition~\cite{nist2023airmf} \\
UNESCO (2021) & Recommandation & Principes~: équité, transparence et
explicabilité, responsabilité, surveillance humaine~\cite{unesco2021ethique} \\
\end{tableaucours}
\vspace{-0.8em}
\begin{center}{\small\itshape Tableau~28.1. Trois références pour auditer
un système d'IA}\end{center}

% ============================================================================
\section{Le Programme d'Audit d'un Modèle}
\label{sec:c28-programme}
% ============================================================================

\subsection{Cinq sous-domaines}

\begin{tableaucours}{cL{2.8cm}L{5.6cm}L{2.6cm}}{%
    \tableauHeader{Code} &
    \tableauHeader{Sous-domaine} &
    \tableauHeader{Contrôles} &
    \tableauHeader{AI RMF}
}
IA.1 & Gouvernance & Inventaire des modèles~; propriétaire métier
désigné~; approbation avant usage & Gouverner \\
IA.2 & Données & Provenance et représentativité documentées~; licéité
du traitement (loi n°~2024/017, Chapitre~\ref{chap:droit}) &
Cartographier \\
IA.3 & Conception et validation & Validation par une équipe indépendante
du concepteur~; motifs de décision communicables à la personne concernée
& Mesurer \\
IA.4 & Biais et équité & Mesure périodique des écarts entre groupes~;
revue des variables corrélées à un groupe protégé & Mesurer \\
IA.5 & Exploitation et dérive & Performance suivie en production~;
réentraînement soumis à la gestion des changements~; recours humain &
Gérer \\
\end{tableaucours}
\vspace{-0.8em}
\begin{center}{\small\itshape Tableau~28.2. Décomposition de l'audit d'un
système d'IA}\end{center}

\subsection{Mesurer un biais, mesurer une dérive}

Le cours GCI définit le \emph{ratio d'impact}~: taux de décisions
favorables du groupe le moins favorisé, divisé par celui du groupe le plus
favorisé. Une pratique issue du droit américain du travail tient pour
suspect un ratio inférieur à~0,8~; ce seuil n'a pas de valeur juridique au
Cameroun, il sert de signal d'alerte. La dérive se mesure en comparant la
performance prévue lors de la validation à la performance constatée en
production.

\begin{boxscenario}{Opération SANAGA. Le modèle de notation}
\textbf{Tests de l'auditrice} sur les 6~240~demandes de microcrédit du
premier semestre~2026 (extraction tracée, Chapitre~\ref{chap:preuve}).

\medskip
\textbf{Biais.} Taux d'acceptation de 42~\% dans les agences du Nord,
contre 70~\% dans celles du Littoral~: ratio d'impact
$= 0{,}42 / 0{,}70 = 0{,}60$, sous le seuil d'alerte. Parmi les variables
du modèle figure l'agence de dépôt, corrélée à la région.

\medskip
\textbf{Dérive.} Le dossier de validation prévoyait un taux de défaut de
6~\%~; il atteint 11~\% sur les crédits accordés depuis janvier. Le modèle
a été réentraîné en décembre~2025 par le prestataire, sans passer par la
gestion des changements de la banque.

\medskip
\textbf{Analyse} (confiance moyenne). L'écart entre régions peut refléter
en partie des différences réelles de risque~; seule une analyse à risque
égal, sans la variable «~agence~», permettra de conclure. La dérive, elle,
est établie, et coïncide avec un réentraînement non contrôlé.
\end{boxscenario}

\begin{boxlecture}{Modèle de notation, SANAGA}
Dix contrôles, deux par sous-domaine~:

\medskip
\begin{tableaucours}{L{7.0cm}c}{%
    \tableauHeader{Contrôle} &
    \tableauHeader{Cote}
}
IA.1 Modèle inscrit à l'inventaire & \OUI{} \\
IA.1 Propriétaire métier désigné & \NON{} \\
IA.2 Provenance des données documentée & \OUI{} \\
IA.2 Licéité du traitement analysée & \NEK{} \\
IA.3 Validation indépendante du concepteur & \NON{} \\
IA.3 Motifs de refus communiqués au client & \NON{} \\
IA.4 Écarts entre régions mesurés périodiquement & \NON{} \\
IA.4 Variables corrélées revues & \NEK{} \\
IA.5 Taux de défaut suivi chaque mois & \OUI{} \\
IA.5 Réentraînement soumis aux changements & \NON{} \\
\end{tableaucours}

\medskip
\textbf{Constatation.} $\mu = 0{,}30$~; $U = 0{,}20$~; $H \approx 1{,}49$~bit.

\medskip
\textbf{Analyse} (confiance moyenne). Non-conformité opaque~: la banque
suit le taux de défaut, mais ne gouverne pas le modèle. Les deux \NEK{}
portent sur des pièces que le prestataire n'a pas remises~; sans droit
d'audit au contrat (Chapitre~\ref{chap:audit-continuite}), la banque ne
peut pas les obtenir seule. Recommandation~: suspendre la variable
«~agence~» en attendant l'analyse à risque égal, et soumettre tout
réentraînement à la gestion des changements.

\medskip
\hfill\lecture{0{,}30}{0{,}20}{1{,}49}
\end{boxlecture}

% ============================================================================
\section{Un Exemple de Régulation de l'IA}
\label{sec:c28-regulation}
% ============================================================================

Le Cameroun n'a pas, à la date de rédaction, de loi propre à l'IA~; les
décisions automatisées qui touchent des personnes relèvent notamment de la
loi n°~2024/017 sur les données personnelles (à vérifier). L'Union
européenne a adopté en 2024 un règlement sur l'intelligence
artificielle~\cite{ue2024aiact}, qui illustre une approche par les
risques~: certaines pratiques sont interdites, les systèmes « à haut
risque » (par exemple l'évaluation de la solvabilité des personnes) sont
soumis à des exigences de gestion des risques, de qualité des données, de
documentation, de transparence et de surveillance humaine.

\begin{boxinfo}
\textbf{Pourquoi l'auditeur camerounais s'y intéresse.} Le modèle de
notation des microcrédits de la BANQUE SANAGA serait, dans l'Union
européenne, un système à haut risque. Sans valeur juridique au Cameroun,
les exigences du règlement fournissent une liste de critères mûrie, que
l'auditeur peut proposer comme référentiel volontaire, en le déclarant
comme tel dans le plan d'audit.
\end{boxinfo}

% ============================================================================
\section{Auditer avec l'IA Générative}
\label{sec:c28-avec}
% ============================================================================

\subsection{Un outil transversal}

L'article A0 propose de traiter l'IA générative comme un \emph{outil
transversal} au service des piliers de l'auditologie, et non comme un
sixième pilier~: elle augmente la capacité de traitement de l'auditeur et
le débit de l'information qui lui parvient, sans changer la nature de ce
qu'il cherche. L'article tient cette position pour révisable~: l'arrivée
d'agents hybrides, humains et artificiels, est un sujet de recherche
ouvert~\cite{minka2026auditologie}.

\begin{tableaucours}{L{3.4cm}L{4.4cm}L{4.4cm}}{%
    \tableauHeader{Usage} &
    \tableauHeader{Apport} &
    \tableauHeader{Limite}
}
Résumer des procédures & Gain de temps en prise de connaissance &
Affirmations inventées~; nuances perdues \\
Proposer des requêtes SQL & Accélère les tests sur données & Requête
fausse mais plausible~; à vérifier sur un cas connu \\
Rédiger des brouillons & Structure, style & Mélange constat et analyse~;
qualification hâtive \\
Classer des tickets, des commentaires & Traite de grands volumes &
Classement non reproductible d'une exécution à l'autre \\
\end{tableaucours}
\vspace{-0.8em}
\begin{center}{\small\itshape Tableau~28.3. Usages de l'IA générative dans
la mission}\end{center}

\begin{boxscenario}{Opération LIZBIZ. Le résumé automatique}
\textbf{Faits.} AUDIX CONSEIL vérifie, procédure en main, 40~affirmations
tirées du résumé produit par l'assistant~: 34 sont exactes, 4 inexactes,
2 invérifiables. L'une des inexactitudes affirme que les droits sont
revus chaque trimestre~; la procédure prévoit une revue annuelle.

\medskip
\textbf{Analyse} (confiance élevée). Un auditeur qui aurait repris le
résumé aurait coté \OUI{} un contrôle qui n'existe pas. Une affirmation de
l'assistant a la valeur d'une déclaration orale~: au mieux, une piste
(Chapitre~\ref{chap:preuve}).
\end{boxscenario}

\begin{boxmetier}{Auditeur SI}
\textbf{Utiliser l'IA générative dans une mission}
\begin{enumerate}
    \item Ne soumettre à un service externe aucune donnée de la mission
          sans l'accord écrit du commanditaire~: le secret de l'auditeur
          (Chapitre~\ref{chap:droit}) s'applique aussi à ce que l'on envoie
          à un outil.
    \item Traiter chaque sortie comme une piste~: toute affirmation est
          vérifiée sur la pièce source avant d'être cotée.
    \item Tester toute requête ou formule proposée sur un cas dont on
          connaît le résultat.
    \item Consigner dans le dossier l'outil, sa version, la demande et
          l'usage fait du résultat.
    \item Ne jamais lever un \NEK{} sur la foi d'une sortie de l'outil.
\end{enumerate}
\end{boxmetier}

\begin{boxmetier}{Gestionnaire de risques}
\textbf{Tenir l'inventaire des modèles d'IA}
\begin{enumerate}
    \item Recenser chaque modèle utilisé dans une décision qui touche des
          clients, des agents ou des comptes.
    \item Pour chacun~: propriétaire métier, fournisseur, données utilisées,
          date de la dernière validation et du dernier réentraînement.
    \item Coter chaque modèle sur l'échelle commune des risques.
    \item Exiger, au contrat de tout fournisseur de modèle, le droit d'audit
          et la remise du dossier de validation.
\end{enumerate}
\end{boxmetier}

% ============================================================================

\separateur

\begin{boxresume}
\textbf{Ce qu'il faut retenir de ce chapitre~:}
\begin{itemize}
    \item \textbf{Dans un modèle, la règle est apprise}~: l'auditeur observe
          les décisions, les données et la surveillance.
    \item \textbf{Cinq sous-domaines}~: gouvernance, données, conception et
          validation, biais et équité, exploitation et dérive.
    \item \textbf{Le ratio d'impact} signale un biais, sans le prouver~; la
          dérive compare le prévu et le constaté~; un réentraînement est un
          changement.
    \item \textbf{L'IA générative est un outil transversal} (A0)~: ses
          sorties sont des pistes, jamais des preuves, et aucune donnée de
          mission ne lui est confiée sans accord.
    \item \textbf{Les cadres de l'IA} (ISO/IEC~42001, NIST AI RMF)
          fournissent les critères~; une régulation de l'IA, comme le
          règlement européen, classe les systèmes par niveau de risque.
\end{itemize}
\end{boxresume}

\bigskip

\begin{boxexo}[title={\ding{45}\quad Exercice~28.1. Rattacher les contrôles \niveaubase}]
Rattachez chaque contrôle à un sous-domaine (IA.1 à IA.5).
\begin{enumerate}[(a)]
    \item Un client refusé peut demander l'examen de son dossier par un
          agent.
    \item Le modèle a été validé par une équipe qui ne l'a pas construit.
    \item Les données d'entraînement couvrent les douze régions.
    \item Chaque modèle a un propriétaire au sein des métiers.
    \item Le taux d'acceptation est comparé chaque trimestre entre régions.
\end{enumerate}
\end{boxexo}

\begin{boxexo}[title={\ding{45}\quad Exercice~28.2. Biais et dérive \niveaumoyen}]
Sur le second semestre, le modèle de la BANQUE SANAGA accepte 48~\% des
demandes dans le Nord et 64~\% dans le Littoral~; le taux de défaut
constaté est de 9~\%.
\begin{enumerate}[(a)]
    \item Calculez le ratio d'impact et comparez-le au premier semestre.
    \item Peut-on conclure que le biais a disparu~? Que faut-il de plus~?
    \item La dérive est-elle corrigée~? Justifiez.
    \item Quelles cotations des dix contrôles changeriez-vous, et pourquoi~?
\end{enumerate}
\end{boxexo}

\begin{boxexo}[title={\ding{45}\quad Exercice~28.3. Une charte d'usage de l'IA générative \niveauexpert}]
AUDIX CONSEIL veut une charte d'usage de l'IA générative pour ses
auditeurs.
\begin{enumerate}[(a)]
    \item Rédigez-en les règles, en vous appuyant sur l'encadré « Utiliser
          l'IA générative dans une mission ».
    \item Quels articles de la loi n°~2010/012 fondent la règle sur les
          données confiées à un service externe~?
    \item Comment rendre reproductible un classement de tickets produit par
          un modèle génératif~?
    \item Discutez la position d'A0~: l'IA générative est-elle un outil, ou
          change-t-elle la nature de l'auditeur~?
\end{enumerate}
\end{boxexo}

\bigskip

\begin{boxexo}[title={\ding{45}\quad Exercice~28.4. Les fonctions de l'AI RMF \niveaubase}]
Rattachez aux fonctions du NIST AI RMF (gouverner, cartographier, mesurer,
gérer)~: (a) l'inventaire des modèles~; (b) le calcul du ratio d'impact~;
(c) la décision de suspendre un modèle qui dérive~; (d) la désignation d'un
propriétaire métier.
\end{boxexo}

\begin{boxexo}[title={\ding{45}\quad Exercice~28.5. La dérive \niveaumoyen}]
Le modèle de la BANQUE SANAGA prévoyait 6~\% de défaut~; on constate
11~\%. Rédigez la constatation et la fiche d'assertion du contrôle « la
performance du modèle est surveillée et les écarts déclenchent une
revue ».
\end{boxexo}

\begin{boxexo}[title={\ding{45}\quad Exercice~28.6. Le réentraînement \niveaumoyen}]
Le modèle a été réentraîné en décembre~2025 hors de la gestion des
changements. Pourquoi est-ce un constat~? À quel chapitre du cours se
rattache le contrôle manquant~?
\end{boxexo}

\begin{boxexo}[title={\ding{45}\quad Exercice~28.7. Une variable suspecte \niveaumoyen}]
La variable « agence » figure parmi les entrées du modèle. Expliquez
comment elle peut produire un biais régional sans que la région soit une
variable du modèle, et ce que l'auditeur demande au fournisseur.
\end{boxexo}

\begin{boxexo}[title={\ding{45}\quad Exercice~28.8. Un résumé à vérifier \niveaumoyen}]
Sur 40~affirmations du résumé automatique des procédures de LiZbiZ, 34 sont
exactes, 4 fausses et 2 invérifiables. Lisez ce résultat avec $\mu$, $U$ et
$H$, et dites ce qu'il implique pour l'usage du résumé dans la mission.
\end{boxexo}

\begin{boxexo}[title={\ding{45}\quad Exercice~28.9. Un référentiel volontaire \niveauexpert}]
Proposez, à partir des exigences d'un système à haut risque du règlement
européen, cinq critères d'audit pour le modèle de notation de la BANQUE
SANAGA, avec leur preuve. Comment le déclarez-vous dans le plan d'audit~?
\end{boxexo}

\begin{boxexo}[title={\ding{45}\quad Exercice~28.10. Suspendre le modèle~? \niveauexpert}]
Le comité d'audit demande s'il faut suspendre le modèle. Rédigez la réponse
de Mme~EKAMBI, en séparant ce que l'audit établit, ce qu'il ne peut pas
établir, et ce qui relève de la décision de la direction.
\end{boxexo}

\bigskip

\begin{boxinfo}
\textbf{Points de contrôle de l'élève.} Je sais~: dire ce que l'IA change
pour l'auditeur~; décomposer l'audit d'un modèle en cinq sous-domaines~;
calculer un ratio d'impact~; repérer une dérive~; fixer les règles d'usage
de l'IA générative dans une mission~; citer deux cadres de l'IA.
\end{boxinfo}

\begin{boxinfo}
\textbf{Pour aller plus loin.}
\begin{itemize}
    \item ISO/IEC~42001:2023~\cite{iso42001_2023}.
    \item NIST, AI RMF~1.0~\cite{nist2023airmf}.
    \item Règlement (UE) 2024/1689~\cite{ue2024aiact}.
    \item UNESCO, Recommandation sur l'éthique de l'IA~\cite{unesco2021ethique}.
\end{itemize}
\end{boxinfo}

\begin{boxinfo}
\textbf{Transition.} Les constats d'un audit peuvent un jour devenir des
pièces d'une procédure judiciaire~: ceux de LiZbiZ intéressent déjà le
ministère public. Le Chapitre~\ref{chap:audit-preuve} montre comment un
constat d'audit devient une preuve.
\end{boxinfo}
